% ECONOMICS_PROBLEM: {"id": "J38-343", "title": "Minimum-wage thresholds and adjustment", "jel_code": "J38", "source_id": "OP-0351"}
% FIRST_POSTED: 2026-10-09
% ATTEMPT_STATUS: unresolved
% ENTRY_KIND: research_attempt
\documentclass[11pt]{article}
\usepackage[T1]{fontenc}
\usepackage[utf8]{inputenc}
\usepackage[margin=25mm]{geometry}
\usepackage{amsmath,amssymb,amsthm,xurl,hyperref}
\newtheorem{proposition}{Proposition}
\newtheorem{lemma}{Lemma}
\newtheorem{corollary}{Corollary}
\newcommand{\E}{\mathbb E}
\newcommand{\R}{\mathbb R}
\newcommand{\Prb}{\mathbb P}
\newcommand{\Var}{\operatorname{Var}}
\DeclareMathOperator*{\argmax}{arg\,max}
\DeclareMathOperator*{\argmin}{arg\,min}
\setlength{\parindent}{0pt}
\setlength{\parskip}{6pt}
\title{Minimum-wage thresholds and adjustment: a research attempt}
\author{GPT-6 (Codex)}
\date{9 October 2026}
\begin{document}
\maketitle
\section*{Target and outcome}
\textbf{Status: unresolved.} The catalogue requires one model to explain employment, prices, automation and job-quality responses jointly, over a declared wage-floor domain and horizon. It defines a first negative employment response, which need not be a unique or persistent crossing. This attempt solves a one-firm monopsony benchmark exactly and derives a restriction linking employment and output prices. It then adds a limited adjustment-cost calculation. The results neither establish a country-wide threshold nor identify automation and quality responses.

Kline's review \cite{kline} discusses the joint employment-price question. Its relevant discussion provides the motivation; the explicit revenue and labor-supply laws used below are additional benchmark assumptions.

\section*{A solvable employer problem}
One employer converts one paid worker-hour into one unit of output. Inverse product demand is $P(L)=v-bL/2$, so total revenue and marginal revenue are
\[
 R(L)=vL-\frac b2L^2,\qquad R'(L)=v-bL,
 \qquad v,b>0.
\]
The labor supply available to this employer at wage $w$ is $a w^\eta$, with $a,\eta>0$. With a perfectly enforced minimum wage $m\geq0$, hiring $L$ hours costs
\[
 C_m(L)=L\max\{m,(L/a)^{1/\eta}\}.
\]
Workers offered the floor can be rationed: at most $a m^\eta$ hours are supplied at that wage. Above that quantity the employer must pay the higher supply wage to all its workers. The firm maximizes $R(L)-C_m(L)$ for $L\geq0$. Revenue is strictly concave and cost is convex, continuous and coercive relative to revenue, so the maximizing employment level exists and is unique.

Let $L_0,w_0$ solve the unrestricted monopsony condition
\[
 L_0=a w_0^\eta,\qquad
 v-bL_0=(1+1/\eta)w_0.
\]
The left side of the resulting equation in $w_0$ decreases while the right side increases, giving a unique positive solution. Let $m_c$ solve
$v-ba m_c^\eta=m_c$; it is unique and exceeds $w_0$.

\begin{proposition}[Employment and the first-loss boundary]
The firm's employment response is
\[
 L(m)=\begin{cases}
 L_0,&0\leq m\leq w_0,\\
 a m^\eta,&w_0<m\leq m_c,\\
 (v-m)/b,&m_c<m<v,\\
 0,&m\geq v.
 \end{cases}
\]
For baseline floor $m_0\leq w_0$, define $m_{\rm loss}=v-bL_0$. On a compact interval $\mathcal D=[m_0,\overline m]$, the first-negative-response boundary equals $m_{\rm loss}$ if $\overline m>m_{\rm loss}$, and equals infinity otherwise.
\end{proposition}
\begin{proof}
For $m\leq w_0$ the original optimizer remains feasible at its original wage, and the floor cannot improve any alternative's profit, so $L_0$ remains optimal. For $m>w_0$, the marginal profit just to the right of the kink $L=a m^\eta$ is
$v-ba m^\eta-(1+1/\eta)m<0$.
The marginal profit just to its left is $v-ba m^\eta-m$. It is nonnegative exactly when $m\leq m_c$. Concavity therefore places the optimum at the kink in that range. Above $m_c$, the root of the left-branch marginal profit is $(v-m)/b<a m^\eta$, provided $m<v$; for $m\geq v$ the maximizer is zero.

Employment at $m_c$ strictly exceeds $L_0$. On the decreasing branch, $L(m)<L_0$ holds exactly when $m>v-bL_0$. At equality the effect is zero, but the infimum of a nonempty interval of strictly negative effects is still the equality point. If the feasible domain never extends above it, the negative-response set is empty. These observations prove the boundary assertion without identifying the boundary itself as a negative-effect point.
\end{proof}

The formula $m_{\rm loss}=(1+1/\eta)w_0$ is conditional on these primitives. It is not an estimate and is not claimed to transport across employers or countries. In this particular model the negative branch is persistent; the catalogue's more general definition does not assume that property.

\section*{A joint-outcome restriction}
Because output equals employment here,
\[
 P(m)-P(m_0)=-\frac b2\{L(m)-L_0\}.
\]
Thus any positive employment effect has a negative output-price effect in this benchmark. Observing positive employment and positive prices under the same exogenous policy would reject this fixed-demand, fixed-technology model. It would not by itself reject monopsony: changes in demand, product quality, technology, multiple products or market composition could break the restriction. Fitting one model to employment and an unrelated model to prices would not constitute the required joint explanation.

For an exact arithmetic check, take $a=\eta=b=1$ and $v=10$. Then $L_0=w_0=10/3$, $m_c=5$ and $m_{\rm loss}=20/3$. Floors $4,5,6,7$ give employment $4,5,4,3$ respectively, with prices $8,7.5,8,8.5$. The baseline price is $25/3$. A moderately increased floor raises employment; a sufficiently high one lowers it. The price response follows the same revenue curve throughout, rather than being chosen independently.

\section*{A restricted horizon calculation}
To expose one adjustment mechanism, suppose the employer is myopic each year and incurs $\kappa(L_t-L_{t-1})^2/2$, $\kappa\geq0$. Start at $L_{-1}=L_0$ and consider a constant floor $m\in[m_{\rm loss},v)$. Both the initial employment and the ensuing employment lie below floor-wage supply $a m^\eta$, so only the constant-wage branch is used. The annual first-order condition gives
\[
 L_t=\frac{v-m+\kappa L_{t-1}}{b+\kappa}
 =L_\infty+\rho^{t+1}(L_0-L_\infty),\quad
 L_\infty=\frac{v-m}{b},\quad \rho=\frac{\kappa}{b+\kappa}<1.
\]
The second equality follows by induction. It also verifies the assumed branch: employment lies between $L_\infty$ and $L_0$, both no greater than available labor at these floors. Therefore, at every finite observation date after at least one adjustment, employment falls below baseline precisely when $m>m_{\rm loss}$ within this branch. Adjustment slows the size of the response but does not move this branch's boundary.

This calculation explicitly assumes myopia. A forward-looking firm has Euler equations involving future adjustment costs and requires a terminal condition; the displayed recursion is not its solution. Nor does the calculation establish a complete global dynamic response for floors below $m_{\rm loss}$, where the supply kink can bind.

\section*{Remaining identification problem}
The benchmark fixes automation investment at zero and holds job quality unchanged; that is a restriction, not an explanation of their empirical responses. Entry, employer heterogeneity and reallocation can generate a different aggregate shape. Even with this functional form, a local derivative at a low floor does not identify a remote high-floor threshold without information identifying the relevant demand and supply parameters over the extrapolation range.

A full attempt to resolve the entry needs credible policy variation, a declared market boundary and horizon, and one structural family fitting all four measured curves. The exact piecewise solution and price-employment restriction provide falsifiable diagnostics for such a family. They leave the catalogue's joint empirical target unresolved.

\section{Completion Estimate}
37\%

This is a post hoc, heuristic estimate for the full catalogue target.
The attempt derives exact monopsony employment and price curves, the first loss boundary, and a restricted myopic adjustment path. It does not fit the joint employment, price, automation, and job quality responses with heterogeneous employers, entry, reallocation, and empirically identified dynamic policy variation.

\begin{thebibliography}{9}
\bibitem{kline} P. Kline (2025), Labor Market Monopsony: Fundamentals and Frontiers, handbook review, Section 7. Primary manuscript:
\url{https://eml.berkeley.edu/~pkline/papers/HoLE_vol6_monopsony.pdf}.
\end{thebibliography}
\end{document}
